\section{Бикомпактное и компактное множества}

\begin{definition}[Бикомпактное множество]
  Подмножество топологического пространства называется \emph{бикомпактным}, если любое его открытое покрытие имеет конечное подпокрытие.
\end{definition}

\begin{definition}[Компактное множество]
 Замкнутое бикомпактное подмножество топологического пространства называется \emph{компактным}. 
\end{definition}

\newpage

\section{Полуметрика и метрика}

\begin{definition}[Полуметрика, метрика]
  Пусть $X$ -- непустое множество. Отображение
  \[
    d : x^2 \rightarrow \R
  \]
  называется \emph{метрикой} на $X$, если для любых элементов $x,y,z \in X$ выполнены следующие условия:
  \begin{enumerate}
    \item $d(x,y) = 0 \iff x = y$.
      \item $d(x,y) = d(y,x)$ (симметричность).
        \item $d(x,z) \leqslant d(x,y) + d(y,z)$ (неравенство треугольника).
  \end{enumerate}
  
  В случае, когда вместо первого условия выполнено более слабое условие
  \[
    d(x,x) = 0,
  \]
  отображение $d$ называется \emph{полуметрикой}.
\end{definition}

\section{Равномерно непрерывное отображение полуметрических пространств}

\begin{definition}[Равномерно непрерывное отображеие полуметричесх пространств]
  Отображение $f$ из полуметрического пространства $(X,d)$ в полуметрическое пространство $(Y,\rho)$ называется \emph{равномерно непрерывным на множестве} $A \subset X$, если выполнены условия:
  \begin{enumerate}
    \item $A \subset \dom f$.
      \item $\forall \epsilon > 0 \ \exists \delta > 0 : \ \forall x_1,x_2 \in A \ d(x_1,x_2) < \delta \implies \rho \big(f(x_1),f(x_2)\big) < \epsilon $.
  \end{enumerate}
\end{definition}

\section{Последовательность Коши}

\begin{definition}[Последовательность Коши]
  Последовательность ${x_n}_{n\in\N}$ элементов полуметрического пространства $(X,d)$ называется \emph{последовательностью Коши}, или \emph{фундаментальной последовательностью}, если
  \[
    \forall \epsilon > 0 \exists n_0 \in \mathbb{N} : \forall n,m \geqslant n_0 \ d(x_n,x_m) < \epsilon.
  \]
\end{definition}

\section{Полнота подмножества полуметрического пространства}

\begin{definition}[Полное подмножество]
  Подмножество полуметрического пространства называется \emph{полным}, если любая его последовательность Коши сходится к некоторому элементу этого подмножества.
\end{definition}

\section{Принцип вложенных шаров}

\begin{theorem}[Принцип вложенных шаров]
  Для того чтобы полуметрическое пространство было полным, необходимо и достаточно, чтобы любая последовательность замкнутых, вложенных друг в друга шаров, радиусы которых стремятся к нулю, имела непустое пересечение.
\end{theorem}

\section{Неподвижная точка, условие Липшица, сжимающее отображение}

\begin{definition}[Неподвижная точка]
  Пусть $f$ -- отображение множества $X$ в себя. Тогда $x \in X$ называется \emph{неподвижной точкой отображения} $f$, если $f(x) = x$.
\end{definition}

\begin{definition}[Условие Липшица]
  Говорят, что отображение $f$ из полуметрического пространства $(X,d)$ в полуметрическое пространство $(Y,\rho )$ удовлетворяет на множестве $A \subset \dom f$ \emph{условию Липшица с константой} $L>0$, если
  \[
    \forall x ', x '' \in A \ \rho \big(f(x '), f(x '')\big) \leqslant L \cdot d (x ', x ''). 
  \]
\end{definition}

\begin{definition}[Сжимающее отображение]
  Отображение из одного полуметрического пространства в другое называется \emph{сжимающим}, если оно является липшицевым с константой $L < 1$. Другими словами, отображение $f$ из полуметрического пространства $(X,d)$ в полуметрическое пространство $(Y,\rho )$ называется \emph{сжимающим}, если для некоторого числа $L \in (0,1)$ выполнено условие
  \[
    \forall x ', x '' \in \dom f \ \rho \big(f(x '), f(x '')\big) \leqslant L \cdot d(x ', x ''). 
  \]
\end{definition}

\section{Теорема Банаха о неподвижной точке}

\begin{theorem}[Банаха о неподвижной точке]
  Каждое сжимающее отображение полного метрического пространства в себя имеет, и притом только одну, неподвижную точку.
\end{theorem}

\section{Линейное отображение}

\begin{definition}[Линейное отображение]
  Пусть $X$ и $Y$ -- векторные пространства над полем $\mathbb{F}$. Отображение $f: X \rightarrow Y$ называется \emph{линейным отображением} или \emph{линейным оператором}, если для любых элементов $x_1,x_2 \in \dom f, \ \lambda , \mu \in \mathbb{F}$
  \[
    \lambda x_1 + \mu x_2 \in \dom f \text{ и } f(\lambda x_1 + \mu x_2) = \lambda f(x_1) + \mu f(x_2).
  \]
\end{definition}

\section{Ядро линейного оператора}

\begin{definition}[Ядро линейного оператора]
  Пусть $X$ и $Y$ -- векторные пространства над полем $\mathbb{F}$. \emph{Ядром линейного отображения} $A$ из $X$ в $Y$ называется прообраз нуля, то есть множество
  \[
    \ker A \coloneq A^{-1} [\mathbb{O}] = \big\{x \in \dom A \ \big| \ Ax = \mathbb{O}\big\}.
  \]
\end{definition}

(должна быть маленькая $\mathbb{O}$)

\section{Критерий взаимной однозначности линейного отображения}

\begin{theorem}[Критерий взаимной однозначности]
  Пусть $X$ и $Y$ -- векторные пространства над полем $\mathbb{F}$. Для того чтобы линейный оператор $A$ из $X$ в $Y$ был взаимно однозначным, необходимо и достаточно, чтобы его ядро состояло из одного нуля: $\ker A = \mathbb{O}$ (тут маленькая о, хз как ее написать). 
\end{theorem}
